\textbf{(CAP-COLUNI 2023 -- Q20)} No painel de um carro, o mostrador do nível de combustível é formado por dois retângulos e dois trapézios retângulos, conforme figura abaixo.

\includegraphics{./figs/2023-Q20.png}

Os dois retângulos possuem dimensões iguais a 0,5 cm e 1,0 cm, um dos trapézios possui bases iguais a 0,5 cm e 0,75 cm, além de altura igual a 1,0 cm. O outro trapézio (aquele que se encontra no topo da figura), possui bases iguais a 0,75 cm e 1,0 cm e 1,0 cm de altura.

Neste mostrador, há um indicador do percentual de ocupação do tanque que marca 100\% quando ele está completamente preenchido e 0\% quando o tanque está completamente vazio. Porém, analisando geometricamente a estrutura do mostrador, vemos que a indicação de 50\% de ocupação do tanque (a 2,0 cm do eixo horizontal) não condiz com a metade da área do polígono que forma o mostrador.

É CORRETO afirmar que a indicação de 50\% de ocupação do tanque deveria estar a uma distância em relação ao eixo horizontal igual a (se necessário use $\sqrt2=1{,}4$ e $\sqrt3=1{,}7$):

\begin{enumerate}
	\item 2,12 cm
	\item 2,25 cm
	\item 2,38 cm
	\item 2,43 cm
\end{enumerate}
